%% This file contains a number of tweaks that are typically applied to the main document.
%% They are not enabled by default, but can be enabled by uncommenting the relevant lines.

%%
%% Inline annotations; for predefined colors, refer to "dvipsnames" in the xcolor package:
%% https://tinyurl.com/overleaf-colors
%%
\newcommand{\red}[1]{{\color{red}#1}}
\newcommand{\todo}[1]{{\color{red}#1}}
\newcommand{\TODO}[1]{\textbf{\color{red}[TODO: #1]}}
%%
%% disable for camera ready / submission by uncommenting these lines  
%%
% \renewcommand{\TODO}[1]{}
% \renewcommand{\todo}[1]{#1}

%%
%% work harder in optimizing text layout. Typically shrinks text by 1/6 of page, enable
%% it at the very end of the writing process, when you are just above the page limit
%%
% \usepackage{microtype}

%%
%% fine-tune paragraph spacing
%%
% \renewcommand{\paragraph}[1]{\vspace{.5em}\noindent\textbf{#1.}}

%%
%% globally adjusts space between figure and caption
%%
% \setlength{\abovecaptionskip}{.5em}


%%
%% Allows "the use of \paper to refer to the project name"
%% with automatic management of space at the end of the word
%%
% \usepackage{xspace}
% \newcommand{\paper}{ProjectName\xspace}

%%
%% Commonly used math definitions
%%
% \DeclareMathOperator*{\argmin}{arg\,min}
% \DeclareMathOperator*{\argmax}{arg\,max}

%%
%% Tigthen underline
%%
% \usepackage{soul}
% \setuldepth{foobar}

\usepackage[dvipsnames]{xcolor}
\usepackage{multirow}
\usepackage{makecell}
% \usepackage{todonotes}
\usepackage{cuted}
\usepackage{tabularx}
\usepackage{chngcntr}
\usepackage{colortbl}
\usepackage{caption}
\usepackage{bm}
\usepackage{listings}
\usepackage{afterpage}

\usepackage[accsupp]{axessibility}  % Improves PDF readability for those with disabilities.

\usepackage{tcolorbox}
\tcbuselibrary{listings, breakable}

\definecolor{lightgray}{gray}{0.95}

\newtcblisting{PromptBox}{
%   breakable,
  colback=lightgray,
  colframe=black!30,
  listing only,
  listing options={
    basicstyle=\ttfamily\footnotesize,
    breaklines=true,
    columns=fullflexible,
    keepspaces=true,
    aboveskip=0pt,
    belowskip=0pt,
    lineskip=-1pt,
  },
  left=3pt,
  right=3pt,
  top=3pt,
  bottom=3pt,
  before=\setlength{\parskip}{0pt},
  before skip=6pt,
}

% \lstset{
%   breaklines=true,
%   breakatwhitespace=false,
%   basicstyle=\ttfamily\footnotesize,
%   columns=fullflexible,
%   backgroundcolor=\color{gray!5},
%   frame=single,
%   framexleftmargin=0pt,
%   % aboveskip=0pt,
%   % belowskip=0pt,
%   lineskip=-1pt,
% }

\newcommand{\firstcolor}{red!30}
\newcommand{\firstcell}[1]{\cellcolor{\firstcolor}#1}
\newcommand{\secondcolor}{orange!30}
\newcommand{\secondcell}[1]{\cellcolor{\secondcolor}#1}
\newcommand{\thirdcolor}{yellow!30}
\newcommand{\thirdcell}[1]{\cellcolor{\thirdcolor}#1}
\newcommand{\graycell}[1]{\cellcolor{gray!10}#1}

\newcommand{\firstgray}{gray!50}
\newcommand{\firstgraycell}[1]{\cellcolor{\firstgray}#1}
\newcommand{\secondgray}{gray!30}
\newcommand{\secondgraycell}[1]{\cellcolor{\secondgray}#1}
\newcommand{\thirdgray}{gray!10}
\newcommand{\thirdgraycell}[1]{\cellcolor{\thirdgray}#1}
